\documentclass[11pt,a4paper]{article}
\usepackage[T1]{fontenc}
\usepackage[utf8]{inputenc}
\usepackage{amsmath,amssymb,amsthm}
\usepackage[margin=1in]{geometry}
\usepackage[colorlinks=true,linkcolor=blue,urlcolor=blue]{hyperref}
\theoremstyle{plain}
\newtheorem{theorem}{Theorem}
\newtheorem{lemma}[theorem]{Lemma}
\newtheorem{proposition}[theorem]{Proposition}
\newtheorem{corollary}[theorem]{Corollary}
\theoremstyle{definition}
\newtheorem*{remark}{Remark}

\title{The prime factorisation of a triple product of gcds,\\
and an off-by-one in the conjecture as posted}
\author{Adrian Perez Fontelles\\ \small Independent researcher}
\date{23 August 2026}
\begin{document}
\maketitle

\begin{abstract}
OEIS A129454 is $a(n)=\prod_{i=1}^{n-1}\prod_{j=1}^{n-1}\prod_{k=1}^{n-1}\gcd(i,j,k)$. In
April 2007 P.~Bala conjectured that
$\operatorname{ord}_{p}a(n)=\sum_{t\ge1}\lfloor n/p^{t}\rfloor^{3}$. We prove the exact
prime factorisation
$\operatorname{ord}_{p}a(n)=\sum_{t\ge1}\lfloor(n-1)/p^{t}\rfloor^{3}$, by the counting
argument that also gives Legendre's formula. The conjecture as literally indexed is false
--- the comment was transplanted from the companion entry A092287, whose products run to
$n$ rather than $n-1$, and the upper limit was not adjusted. The smallest counterexample is
$n=4$, $p=2$: the posted formula gives exponent 9, the truth is 1, and indeed $a(4)=6$.
With $n$ replaced by $n-1$ throughout, the conjecture is exactly right, and that is the
reading the entry's own definition forces.
\end{abstract}

\section{The sequence and the conjecture}

OEIS A129454, contributed by Peter Bala on 16 April 2007, has offset 0 and
\[
  a(n)=\prod_{i=1}^{n-1}\prod_{j=1}^{n-1}\prod_{k=1}^{n-1}\gcd(i,j,k)\quad(n>2),
  \qquad a(n)=1\ (n\le2),
\]
with $a(0),\dots,a(7)=1,1,1,2,6,1536,7680,8806025134080$. The entry carries

\begin{quote}
\textbf{Conjecture:} Let $p$ be a prime and let $\operatorname{ord}_{p}(n,p)$ denote the
largest power of $p$ which divides $n$. \dots\ Then we conjecture that the prime
factorization of $a(n)$ is given by $\operatorname{ord}_{p}(a(n),p)=(\lfloor
n/p\rfloor)^{3}+(\lfloor n/p^{2}\rfloor)^{3}+(\lfloor n/p^{3}\rfloor)^{3}+\cdots$. Compare
with the comments in A092287.
\end{quote}

As of the ``Last modified'' line on the live entry (23 August 2026) this is still recorded
as an unproven conjecture, with no proof in the entry's links or programs.

The pointer to A092287 explains the shape of the statement and also its defect. In A092287
the products run over $1\le j,k\le n$ and the correct exponent is
$\sum_{t}\lfloor n/p^{t}\rfloor^{2}$; here the products run over $1\le i,j,k\le n-1$, but
the comment reuses $\lfloor n/p^{t}\rfloor$ unchanged. As stated it is therefore false.
Section~3 proves the corrected version and Section~4 records the counterexample.

Throughout $p$ is a prime and $v_{p}(N)$ is the exponent of $p$ in $N$.

\section{The counting lemma}

\begin{lemma}[layer cake]\label{lem:layer}
For positive integers $x_{1},\dots,x_{d}$,
\[
  v_{p}\bigl(\gcd(x_{1},\dots,x_{d})\bigr)=\min_{1\le t\le d}v_{p}(x_{t})
  =\#\bigl\{\,s\ge1:p^{s}\mid x_{t}\text{ for every }t\,\bigr\}.
\]
\end{lemma}

\begin{proof}
The first equality is the description of gcd prime by prime. For the second, $p^{s}$
divides every $x_{t}$ exactly when $s\le v_{p}(x_{t})$ for every $t$, that is when
$1\le s\le\min_{t}v_{p}(x_{t})$.
\end{proof}

\section{The theorem}

\begin{theorem}\label{thm:main}
Let $N\ge1$, $d\ge1$ and put
$G_{d}(N)=\prod_{x_{1}=1}^{N}\cdots\prod_{x_{d}=1}^{N}\gcd(x_{1},\dots,x_{d})$. Then for
every prime $p$,
\[
  v_{p}\bigl(G_{d}(N)\bigr)=\sum_{s\ge1}\left\lfloor\frac{N}{p^{s}}\right\rfloor^{d}.
\]
\end{theorem}

\begin{proof}
By additivity of $v_{p}$ and Lemma~\ref{lem:layer},
\[
  v_{p}\bigl(G_{d}(N)\bigr)=\sum_{x_{1},\dots,x_{d}\le N}
  \#\{s\ge1:p^{s}\mid x_{1},\dots,p^{s}\mid x_{d}\}
  =\sum_{s\ge1}\#\bigl\{(x_{1},\dots,x_{d})\in[1,N]^{d}:p^{s}\mid x_{t}\ \forall t\bigr\},
\]
a finite rearrangement of nonnegative terms. The $d$ divisibility conditions are
independent, so the inner count is
$\bigl(\#\{x\le N:p^{s}\mid x\}\bigr)^{d}=\lfloor N/p^{s}\rfloor^{d}$.
\end{proof}

\begin{corollary}\label{cor:main}
For A129454 and every prime $p$,
\[
  \operatorname{ord}_{p}a(n)=\sum_{s\ge1}\left\lfloor\frac{n-1}{p^{s}}\right\rfloor^{3}.
\]
This is Bala's conjecture with $n$ replaced by $n-1$, and it is the statement forced by the
entry's own definition of $a(n)$ as a product over $1\le i,j,k\le n-1$.
\end{corollary}

\begin{proof}
Apply Theorem~\ref{thm:main} with $d=3$ and $N=n-1$. For $n\le2$ the product is empty and
both sides are 0.
\end{proof}

\section{The conjecture as posted is false}

\begin{proposition}\label{prop:false}
The formula $\operatorname{ord}_{p}a(n)=\sum_{s\ge1}\lfloor n/p^{s}\rfloor^{3}$ fails
whenever $p\mid n$ and $n\ge4$. The smallest failure is $n=4$, $p=2$.
\end{proposition}

\begin{proof}
By Corollary~\ref{cor:main} the truth is $\sum_{s}\lfloor(n-1)/p^{s}\rfloor^{3}$, and
$\lfloor n/p^{s}\rfloor=\lfloor(n-1)/p^{s}\rfloor$ unless $p^{s}\mid n$, in which case
$\lfloor n/p^{s}\rfloor=\lfloor(n-1)/p^{s}\rfloor+1$. So if $p\mid n$ the two expressions
differ, the posted one being strictly larger.

Explicitly at $n=4$, $p=2$: the only triples $(i,j,k)\in[1,3]^{3}$ with $\gcd(i,j,k)>1$ are
$(2,2,2)$ and $(3,3,3)$, so $a(4)=2\cdot3=6$ and $v_{2}(a(4))=1$.
Corollary~\ref{cor:main} gives $\lfloor3/2\rfloor^{3}+\lfloor3/4\rfloor^{3}=1+0=1$,
correct. The posted formula gives $\lfloor4/2\rfloor^{3}+\lfloor4/4\rfloor^{3}=8+1=9$.
Since $2^{9}\nmid6$, the posted formula is false at $n=4$. No smaller $n$ fails: for
$n\le3$ every relevant $\lfloor n/p^{s}\rfloor$ agrees with $\lfloor(n-1)/p^{s}\rfloor$.
\end{proof}

\begin{remark}[what to record on the entry]
The right correction is to replace $n$ by $n-1$ in the comment, not to change the
definition: the definition, the DATA and the Mathematica, Magma and SageMath programs on
the entry all agree on products over $1\le i,j,k\le n-1$. It is only the comment,
transplanted from A092287, that is out of step.
\end{remark}

\begin{remark}[honest assessment]
Theorem~\ref{thm:main} is a three-line count and is the same argument that proves
Legendre's formula; its $d=2$ case is the rectangular statement that A092287 records. What
is worth recording here is not the proof but the off-by-one: the statement people have been
checking is not the statement that is true.
\end{remark}

\section{Verification}

A verification script, written separately from this note, checks every claim and prints
every number quoted. It (i) recomputes $a(0),\dots,a(11)$ of A129454 directly from the
triple product and matches the published DATA exactly; (ii) compares $v_{p}(a(n))$ against
$\sum_{s}\lfloor(n-1)/p^{s}\rfloor^{3}$ for every $2\le n\le25$ and every prime $p<n$, with
0 mismatches; (iii) compares it against the posted formula
$\sum_{s}\lfloor n/p^{s}\rfloor^{3}$ over the same range, finding 25 mismatches; and
(iv) confirms that the smallest failure is $n=4$, $p=2$, where the true exponent is 1 and
the posted formula claims 9, with $a(4)=6=2\cdot3$.

\begin{thebibliography}{9}
\bibitem{oeis} N.~J.~A. Sloane et al., \emph{The On-Line Encyclopedia of Integer
Sequences}, \url{https://oeis.org/A129454}, entry A129454, consulted 23 August 2026.
\bibitem{oeis2} ibid., entry A092287, \url{https://oeis.org/A092287}.
\bibitem{legendre} A.~M. Legendre, \emph{Th\'eorie des Nombres}, 3rd ed., Firmin Didot,
Paris, 1830.
\end{thebibliography}
\end{document}
